\section{Prepared cycles and the transfer theorem}\label{sec:experiments}
\subsection{Raw causal data and the reset protocol}
Let $\Theta\ne\varnothing$ be a set. A realized $\theta$ is fixed for the entire execution. In each cycle a standard Borel hidden state is prepared, and positive normalized causal kernels produce the next hidden state, a report $Y_t$, an audit mark $Z_t$, and a visible end-of-cycle flag. Positive means nonnegative. Hidden state may contain the necessary physical history. Audit marks are not feedback to the observer. They are actual jointly specified marks of the performed kernels; in a quantum realization, any audit back-action is included in the instrument. We do not posit a joint path of unperformed incompatible audits.

There is a nonempty finite action set $A$ at every ordinary data call, with all its actions legal. The compulsory preparation call and the end flag are prescribed interfaces, not extra decisions. A cycle has an integer length $\tau\ge1$ counting \emph{all} its raw calls, including preparation. The end flag makes $\tau$ a stopping time of visible within-cycle history. At the next preparation, the physical system is freshly prepared with the same $\theta$ and the observer is overwritten by state $1$. No private seed, counter, or record may survive this overwrite. Distinct cycles use independent physical innovations and fresh controller randomization, conditional on $\theta$. Reset calls have a prescribed bounded score; the predictive examples below give them score zero, but still charge their acquisition and time costs.

Reports for action $a$ lie in a standard Borel space $Y_a$. For every model, finite within-cycle intervention word $\mathbf a=(a_1,\ldots,a_t)$, and its prescribed preparation positions, assume
\begin{equation}\label{eq:product}
 P_{\theta,t}^{\mathbf a}\ll \sigma_{a_1}\otimes\cdots\otimes\sigma_{a_t},
\end{equation}
where each $\sigma_a$ is one probability measure common to the family. The reports include the visible flags. A stopped prefix is padded by a fixed cemetery report if necessary, and each reference includes the required atoms. Only finite-history product domination is asserted. Moving or zero densities are allowed. One-time marginal domination is not sufficient.

Decisions lie in a compact metrizable space $D$, and $0\le\ell(d,z)\le1$ is measurable and continuous in $d$ for every $z$. The audit law is part of the raw kernel and does not depend on a private controller seed except through declared actions; $d$ is a readout, not an unrecorded intervention changing the next state. Decision-dependent physical audits can instead be placed in the action interface when that interface satisfies the stated finite-action assumptions. These distinctions make the cycle reward a well-defined random variable.

\begin{definition}[One reused observer]\label{def:program}
A program $\gam\in\Ga$ consists of action rows $\alpha_i\in\Delta(A)$, Borel report rows $k_{ia}:Y_a\to\Delta([M])$, and readouts $d_i\in D$, $i\in[M]=\{1,\ldots,M\}$. At an ordinary call it draws $a$ from $\alpha_i$, receives $y$, draws $j$ from $k_{ia}(y)$, and issues $d_j$ for that call's audit. It then retains $j$, except that the prescribed boundary overwrite returns it to $1$. The same rows are used at all times, cycles, and parameters. Sampling is fresh and conditionally independent. A compulsory preparation is processed by its prescribed overwrite, not by an extra phase-indexed row. All persistent observer registers are included in $[M]$.
\end{definition}
A Borel row need not be computable. A finite executable program will be supplied under the stronger hypotheses of Section~\ref{sec:effective}. The exact program class is not enlarged to history-dependent or phase-dependent policies when a value is compared.

Write $R=\sum_{t=1}^{\tau}\ell_t$ for one cycle's reward, where $\ell_t=\ell(d_{I_t},Z_t)$ at data calls and is the prescribed score at preparation. Thus $0\le R\le\tau$. Put
\begin{equation}\label{eq:cycle}
 \mu_\theta(\gam)=\E_\theta^\gam\tau,\qquad
 r_\theta(\gam)=\E_\theta^\gam R,\qquad
 g_\theta(\gam)=r_\theta(\gam)/\mu_\theta(\gam).
\end{equation}
The actual renewal occupation measure of any within-cycle measurable state $S_t$ is
\begin{equation}\label{eq:occupation}
 \overline\nu_\theta^\gam(B)
 =\frac1{\mu_\theta(\gam)}\E_\theta^\gam
       \sum_{t=1}^{\tau}1_{\{S_t\in B\}}.
\end{equation}
It is a probability measure if every raw call is represented, and a subprobability measure when only scored data calls are retained. It is not generally the preparation law.

\subsection{Return envelope and common values}
Assume there is a decreasing function $a:\N_0\to[0,\infty)$ with
\begin{equation}\label{eq:UI}
 a(0)<\infty,\quad a(T)\longrightarrow0,\qquad
 \sup_{\theta\in\Theta,\gam\in\Ga}\E_\theta^\gam(\tau-T)_+\le a(T).
\end{equation}
This is a certificate for the \emph{actual} controlled cycle law. A bounded mean alone does not imply it. It is imposed on all legal programs before optimization, including the finite tables later used for approximation.

Let $C_N=\sum_{t=1}^N\ell_t$ along the concatenated cycles, starting at a preparation. Define
\begin{align}\label{eq:values}
 J_{\theta,N}(\gam)&=\E_\theta^\gam C_N/N,&
 J_{\theta,\beta}(\gam)&=(1-\beta)\sum_{t\ge1}\beta^{t-1}\E_\theta^\gam\ell_t,\\
 V_N^\Theta(M)&=\inf_{\gam\in\Ga}\sup_\theta J_{\theta,N}(\gam),&
 V_\beta^\Theta(M)&=\inf_{\gam\in\Ga}\sup_\theta J_{\theta,\beta}(\gam),\notag\\
 V_\av^\Theta(M)&=\inf_{\gam\in\Ga}\sup_\theta g_\theta(\gam).&&\notag
\end{align}
The discounted weights sum to one. Scores occur after the report-to-label update. The integer $N$ counts real raw calls, including reset calls, but is an external evaluation budget: the deployed rows do not receive $N$ or a clock. An exact internal deadline is a different protocol.

Set
\begin{equation}\label{eq:moduli}
 b_N=\frac1N\sum_{j=0}^N a(j),\qquad
 c_\beta=(1-\beta)^2\sum_{N\ge1}N\beta^{N-1}b_N.
\end{equation}
Both tend to zero in their respective limits.

\begin{theorem}[Causal renewal transfer]\label{thm:renewal}
Under the prepared-cycle interface, \eqref{eq:product}, and \eqref{eq:UI}, the following hold for every finite $M$.

The average risk $V_\av^\Theta(M)$ is attained by one Borel program. For each fixed model and program, $C_N/N\to g_\theta(\gam)$ almost surely and in mean. Moreover,
\begin{equation}\label{eq:finite-model}
 V_\av^\Theta(M)=
 \sup_{\varnothing\ne F\subset\Theta,\ |F|<\infty}V_\av^F(M).
\end{equation}
A strict target below the full value is therefore excluded by a finite subfamily of fixed models.

Uniformly in every $\theta$ and $\gam$,
\begin{equation}\label{eq:uniform}
 |J_{\theta,N}(\gam)-g_\theta(\gam)|\le b_N,\qquad
 |J_{\theta,\beta}(\gam)-g_\theta(\gam)|\le c_\beta.
\end{equation}
The same inequalities compare the three optimized values. Every cluster point, in the controller topology described below, of asymptotically optimal finite-acquisition or vanishing-discount programs is average-optimal. Exact optimal programs at finite $N$ and $0<\beta<1$ also exist.

If, in addition, for some $0<p\le1$,
\begin{equation}\label{eq:moment}
 \sup_{\theta,\gam}\E_\theta^\gam\tau^{1+p}\le K,
\end{equation}
then $b_N$ and $c_\beta$ in \eqref{eq:uniform} may be replaced by
\begin{equation}\label{eq:moment-rates}
 K\frac{(N+1)^{1-p}}{N}\le2K N^{-p},\qquad
 2^{1-p}K(1-\beta)^p,
\end{equation}
respectively. Replacing a physical realization by another with the same finite report, survival and payoff responses leaves all these values unchanged, provided the same reset and resource protocol is preserved.
\end{theorem}
There is no minimax interchange in \eqref{eq:finite-model}; the parameter is not resampled or adversarially changed at each call. The theorem gives no cardinality bound on the witnessing finite family. For $0\le\ell\le H$, all risk-error bounds acquire a factor $H$, including a discounted prefix tail $H\beta^T$.

\subsection{Predictive equivalence and its measurable realization}
Fix a preparation, or a specified Bayesian preparation of an unknown parameter. Take a countable determining family of executable future tests, containing the visible interface and closed under the finite continuations needed to determine subsequent tests. Evaluation gives a Borel map
\begin{equation}\label{eq:quotient}
 q(h)=(\E[f_j\mid h])_{j\ge1}\in[0,1]^{\N}.
\end{equation}
The conditional versions are fixed on the actual prepared history space. No model-dependent $q$ is provided to the robust controller as a hidden register.

\begin{proposition}[Predictive realization]\label{prop:quotient}
The fibers of \eqref{eq:quotient} are exactly future-test equivalence. Every sufficient statistic through which all test expectations factor also determines $q$. If the evaluation image is Borel, and the one-step future-test closure includes the joint next-report/test laws, the image is a standard Borel predictive realization: report laws and a Borel update descend to it, up to their actual null sets. Without the Borel-image condition, only the evaluation statistic and its stated minimality are asserted.
\end{proposition}
\begin{proof}
Countable determining tests and the monotone-class principle identify the future laws, proving the fiber assertion. Coordinatewise factorization proves minimality. Suppose the image is Borel. A Borel map constant on the fibers of a Borel surjection between standard Borel spaces has a Borel factor: the image of a saturated Borel set and of its complement are disjoint analytic sets whose union is the target, and hence both are Borel. Apply this observation to the probabilities of a countable generating algebra of next-report/future-test events. It gives the descended joint kernels. Standard Borel disintegration then gives conditional next-test expectations as Borel functions of $(q,a,y)$. Their countably many versions agree with the original updated evaluations almost surely. Complete the update on the remaining Borel null fibers by a fixed image point. This proves realization without postulating a manifold or a Fisher matrix. The underlying disintegration and version facts are standard \cite{Kallenberg,BertsekasShreve}.
\end{proof}
